\BeleskeID{09_3-s046}
% element: 09_3-s046:e05
% slide-id: 09_3-s046
Kod logičkih kola sa bipolarnim tranzistorima istovremeno se koče tranzistori u PUN i PDN mreži. U prikazanoj CMOS realizaciji dodatni tranzistori između PUN/PDN i izlaza prekidaju ili omogućavaju izlazne strujne putanje. Time se ostvaruju visoka impedansa i normalan rad. Kod prikazane realizacije oba dodatna tranzistora su nMOS. Za $E=0$ su zakočeni, a za $E=1$ omogućavaju provođenje. Koja grana zaista provodi zavisi od stanja PUN i PDN.

Postoje i kašnjenja ulaska u visoku impedansu iz nule ili jedinice, $t_{pLZ}$ i $t_{pHZ}$, kao i povratka iz visoke impedanse u aktivnu nulu ili jedinicu, $t_{pZL}$ i $t_{pZH}$. Dodatni tranzistori utiču na ova kašnjenja.

Visoka impedansa $Z$ opisuje električno stanje izlaza, a ne treću Bulovu vrednost. Na zajedničkoj liniji samo jedan trostatički izlaz sme u datom trenutku aktivno nametati logički nivo.
